\section{Introduction}
\label{sec:introduction}

Modern machine learning systems have achieved remarkable performance across diverse analytical benchmarks. However, the most accurate predictors---such as deep neural networks or ensemble tree architectures---even domain experts struggle to interpret, creating a tension between accuracy and interpretability.

In critical applications such as clinical oncology, autonomous driving, and financial risk assessment, high predictive accuracy alone is insufficient. Stakeholders require transparent, verifiable rationales before decisions can be safely operationalized~\cite{BARREDOARRIETA202082}. Explainable Artificial Intelligence (XAI) has emerged to bridge this fundamental tension between predictive capacity and human interpretability. Among post-hoc local explanation frameworks, LIME~\cite{ribeiro-etal-2016-trust} and SHAP~\cite{NIPS2017_8a20a862} have become standard industry and research benchmarks. Despite their ubiquity, practitioner reliance on these tools often presumes that high predictive accuracy naturally guarantees fidelity and reliability in the accompanying explanations.

\subsection{Research Question}
\label{sec:guiding-question}

Motivated by this gap, this seminar investigates the central question: \textit{``Does a highly accurate prediction necessarily come with a reliable explanation?''} 

To answer this question, we systematically examine LIME and SHAP across both theoretical foundations and empirical behavior:
\begin{itemize}[leftmargin=1.2em, itemsep=1.5pt]
    \item \textbf{Conceptual Foundations:} Clarifying where post-hoc local methods stand within the broader XAI taxonomy, distinguishing intrinsic from post-hoc models and local from global scopes.
    \item \textbf{Theoretical Formulations \& Worked Examples:} Walking through the exact mathematical objectives and game-theoretic foundations of both methods, accompanied by step-by-step numerical calculations on an illustrative example.
    \item \textbf{Empirical Demonstration:} Evaluating explanations produced for a Random Forest classifier on the Breast Cancer Wisconsin dataset across confident, borderline, and conflicting prediction scenarios.
    \item \textbf{Faithfulness \& Limitations:} Analyzing explanation stability across random seeds, runtime trade-offs, empirical faithfulness via feature deletion, and why post-hoc attributions must not be conflated with causal effects.
\end{itemize}


